% chapters/16/contoh.tex
\section{Dua Contoh Terapan}
\label{sec:contoh-bab-16}

\subsection{Contoh Pertama: Menghitung Jangkauan Lompatan}

\begin{contohbox}
Kelompok Anda menetapkan $v_0 = 620$ piksel per detik dan $g = 1700$ piksel
per detik kuadrat pada levelnya. Petaknya berukuran 32 piksel. Berapa tinggi
lompatan tertinggi dan jarak mendatar terjauh pada ketinggian yang sama?
\end{contohbox}

Tinggi lompatan tertinggi dihitung dari hukum gerak:

\[
h_{\text{maks}} = \frac{v_0^2}{2g} = \frac{620^2}{2 \times 1700}
= \frac{384\,400}{3\,400} = 113{,}06 \text{ piksel},
\]

yaitu $113{,}06 \div 32 = 3{,}53$ petak. Lama di udara dari tanah sampai
kembali ke tanah adalah $t = 2v_0 / g = 2 \times 620 \div 1700 = 0{,}729$
detik. Dengan kecepatan jalan 300 piksel per detik, jarak mendatar terjauhnya

\[
x_{\text{maks}} = 300 \times 0{,}729 = 218{,}8 \text{ piksel},
\]

yaitu $218{,}8 \div 32 = 6{,}84$ petak.

Angka itu memutuskan beberapa hal sekaligus. Celah selebar tujuh petak pada
level Anda tidak dapat dilompati, sedangkan celah enam petak hanya menyisakan
kelonggaran 0,84 petak, dan itu terlalu sempit untuk pemain baru. Rancangan
yang aman adalah celah lima petak, dengan kelonggaran sekitar 1,8 petak.

Kedua angka itu hanya bergantung pada tiga besaran:
kecepatan lompatan, gravitasi, dan kecepatan jalan. Mengubah salah satunya
mengubah seluruh jangkauan level, dan itulah sebabnya besarannya perlu
ditetapkan lebih dahulu, sebagaimana dijelaskan pada Bab~VI.

\subsection{Contoh Kedua: Menelusuri Temuan \istilah{Playtest} yang Tidak Dapat Diulang}

\begin{contohbox}
Catatan penguji pada lembar \istilah{playtest} berbunyi: ``bunyinya kadang
tidak keluar, saya tidak tahu kapan.'' Anda mencoba mengulanginya selama
sepuluh menit dan tidak berhasil. Apakah catatan itu tidak berguna?
\end{contohbox}

Catatan itu berguna, tetapi belum cukup untuk dikerjakan. Yang perlu dilakukan
adalah mengubahnya menjadi catatan yang dapat diulang, dan caranya bukan
menunggu sampai kejadiannya muncul lagi.

\begin{langkah}
  \item \textbf{Catat waktu dan tempatnya.} Bukalah kembali lembar
        \istilah{playtest} itu, bagian waktu penyelesaian dan titik berhenti.
        Keduanya sudah terisi, dan keduanya menyempitkan pencarian: kejadian
        itu muncul pada menit kedua, ketika penguji sedang menuruni tangga.
  \item \textbf{Sempitkan kemungkinannya.} Berdasarkan waktu dan tempatnya,
        susun dugaan: bunyinya berkaitan dengan pendaratan, bukan dengan
        lompatan. Buatlah daftar tindakan yang perlu dicoba, yaitu melompat
        lalu mendarat, mendarat tanpa melompat, dan melompat berulang.
  \item \textbf{Coba setiap dugaan secara terpisah.} Ternyata bunyinya hilang
        ketika pemain melompat kurang dari setengah detik sesudah mendarat.
        Dugaan itu terbukti, dan sekarang catatannya dapat diulang.
  \item \textbf{Perbaiki, lalu periksa ulang.} Setelah bunyinya tidak lagi
        saling memotong, jalankan kembali ketiga tindakan itu untuk memastikan
        temuan itu benar-benar hilang.
\end{langkah}

Tiga hal dapat dipetik dari contoh itu. Pertama, catatan yang tampak tidak
berguna biasanya memuat keterangan yang cukup untuk disempitkan, dan yang
kurang hanyalah waktu dan tempatnya. Kedua, menyempitkan dugaan menjadi
beberapa tindakan yang dapat dicoba jauh lebih cepat daripada mencoba
mengulanginya secara acak. Ketiga, temuan itu baru boleh dipindahkan ke
\istilah{log bug} setelah ada urutan tindakan yang dapat diulang, sebagaimana
dibahas pada Bagian~\ref{sec:log-bug}.
